\section*{Glossary}

\begin{description}[leftmargin=2.5cm, labelindent=0cm]

\item[SiPM] Silicon photomultiplier
\item[PDE] Photo-Detection Efficiency
\item[SPAD] Single Photo-Avalanche Diode
\item[MC] Monochromator

\item $\lambda$ - photon wavelength
\item $\theta$ - incident photon angle in medium
\item $V$ - SiPM operating over-voltage

\item $\bar{n_j} = n+ik$ - the complex refractive index, $k$ being the extinction coefficient, for medium j
\item $\mu = \frac{4\pi k}{\lambda}$ - the attenuation coefficient, dependent on the extinction coefficient $k$
\item $\theta_{2, 3}$ - the photon's angle of incidence within the silicon dioxide (2) and bulk silicon (3), derived from

\item[$\tox$] Average thickness of the surface silicon dioxide (\ce{SiO2}) passivation layer.

\item[$dp^*$] Effective depth below which free electrons are collected and reach the high-field region.

\item[$X_{PN}$] Center of the p–n junction and location of the high-field region; the boundary where avalanches are either electron- or hole-driven.

\item[$dw^*$] Depth above which holes are collected and reach the high-field region.

\item[$P_e(V), P_h(V)$] Probability that a collected electron or hole triggers an avalanche. Parameterized as: $P_e(V; A_e, V_e) = A_e\left(1 - e^{-V / V_e}\right)$

\item[$FF$] Fraction of the photo-sensitive surface area.

\item[$\bm{J}$] The set of device-specific parameters listed above.

\item[$W_p, W_n$] Probability of photoabsorption within the effective p and n regions, respectively.

\item[$iPDE$] Internal detection efficiency — the probability that a photon transmitted into the silicon bulk is detected.


\end{description}
